% ECONOMICS_PROBLEM: {"id": "O33-298", "title": "AI and aggregate productivity", "jel_code": "O33", "source_id": "OP-0263"}
% FIRST_POSTED: 2026-10-09
% ATTEMPT_STATUS: unresolved
% ENTRY_KIND: research_attempt
\documentclass[11pt,a4paper]{article}
\usepackage[T1]{fontenc}
\usepackage[utf8]{inputenc}
\usepackage{lmodern,amsmath,amssymb,amsthm,mathtools}
\usepackage[margin=25mm]{geometry}
\usepackage{microtype,enumitem,hyperref}
\hypersetup{colorlinks=true,urlcolor=blue,linkcolor=blue}
\setlength{\parindent}{0pt}
\setlength{\parskip}{4pt}
\newtheorem{theorem}{Theorem}[section]
\newtheorem{proposition}[theorem]{Proposition}
\newtheorem{lemma}[theorem]{Lemma}
\newtheorem{corollary}[theorem]{Corollary}
\theoremstyle{definition}
\newtheorem{definition}[theorem]{Definition}
\theoremstyle{remark}
\newtheorem{remark}[theorem]{Remark}
\newcommand{\R}{\mathbb{R}}
\newcommand{\E}{\mathbb{E}}
\newcommand{\Prob}{\mathbb{P}}
\newcommand{\ind}{\mathbf{1}}
\DeclareMathOperator{\Var}{Var}
\DeclareMathOperator{\Cov}{Cov}
\DeclareMathOperator{\diag}{diag}
\DeclareMathOperator*{\argmax}{arg\,max}
\DeclareMathOperator*{\argmin}{arg\,min}
\title{AI productivity with task bottlenecks, deployment costs and capacity limits}
\author{GPT 6 Astra Ultra}
\date{9 October 2026}
\begin{document}
\maketitle
\textbf{Status: Unresolved.} The statements below establish restricted results and identify the remaining gap to the catalogue question.

\section{Question and scope}
The target is a finite-horizon productivity contrast for a declared AI capability, including complement and compute constraints. We solve a static endpoint aggregation problem with endogenous allocation of task production. It gives exact resource-adjusted gains, a bottleneck bound, and an explicit capacity-constrained allocation. The transition, entry and causal identification of the primitives remain unresolved.

\section{Task production and a common resource}
There are $J$ tasks with weights $\omega_j>0$, $\sum_j\omega_j=1$, and substitution elasticity $\sigma>0$. For $\sigma\ne1$ define
\[
 Q(q)=\left[\sum_j\omega_j(q_j/\omega_j)^{(\sigma-1)/\sigma}\right]^{\sigma/(\sigma-1)}.
\]
At $\sigma=1$ use $Q(q)=\prod_j(q_j/\omega_j)^{\omega_j}$. Task $j$ costs $c_j>0$ units of a common productive resource per unit of output. Those costs include the declared human, compute and energy requirements in the common unit. Deploying the capability uses an additional fixed $F\in[0,R)$ units of complementary resources out of the total endowment $R>0$. Thus feasibility is $\sum_jc_jq_j\leq R-F$. This is a one-factor economy with fixed tasks and competitive intermediate production, not an unrestricted growth model.

\begin{theorem}[Exact resource-adjusted aggregate output]
Without task capacity bounds, maximal output is $Q^*=(R-F)/P(c)$, where
\[
 P(c)=\begin{cases}
 (\sum_j\omega_jc_j^{1-\sigma})^{1/(1-\sigma)},&\sigma\ne1,\\
 \prod_jc_j^{\omega_j},&\sigma=1.
 \end{cases}
\]
The maximizing allocation is $q_j=Q^*\omega_j(c_j/P)^{-\sigma}$. If baseline costs are $c^0$ with zero fixed deployment cost and capability costs are $c^1$, the output ratio is
\[
 G=(1-F/R)\frac{P(c^0)}{P(c^1)}.\tag{1}
\]
Optional deployment is beneficial in this model exactly when $F/R\leq1-P(c^1)/P(c^0)$, and its output ratio is $\max\{1,G\}$.
\end{theorem}
\begin{proof}
The task aggregator is increasing, concave and homogeneous of degree one. At a positive optimum its derivative is $Q^{1/\sigma}(\omega_j/q_j)^{1/\sigma}$, including the Cobb--Douglas limit. Equalizing derivative per unit of cost yields $q_j=K\omega_jc_j^{-\sigma}$ for a common positive $K$. Substituting this form into the aggregator and the binding resource constraint gives $Q^*=(R-F)/P$ and the stated allocation. Concavity makes these first-order conditions sufficient; positivity and strict task-substitution curvature give the interior optimizer. Taking the ratio across regimes proves (1). Choosing between the feasible baseline and deployment regimes proves the last assertions.
\end{proof}

The differentiable cost index also gives an exact local aggregation weight:
\[
 \frac{\partial\log P}{\partial\log c_j}
 =\frac{\omega_jc_j^{1-\sigma}}{\sum_k\omega_kc_k^{1-\sigma}},
\]
with weight $\omega_j$ at $\sigma=1$. Thus experimental task gains aggregate with endogenous expenditure shares and the fixed-cost term, not necessarily their raw task frequencies.

\section{A capability that improves only some tasks}
Normalize all baseline costs to one. Suppose an AI-eligible set has total weight $a\in(0,1)$, its unit costs fall to $1/g$ with $g\geq1$, and all other costs stay one. This defines precisely the capability being compared.

\begin{proposition}[Bottleneck ceiling]
For $\sigma\ne1$, forced deployment has output ratio
\[
 G(g)=(1-F/R)\bigl[(1-a)+a g^{\sigma-1}\bigr]^{1/(\sigma-1)}.
\]
For $0<\sigma<1$,
\[
 \lim_{g\to\infty}G(g)=(1-F/R)(1-a)^{-1/(1-\sigma)}<\infty.
\]
At $\sigma=1$, $G(g)=(1-F/R)g^a$, while for $\sigma>1$ the ratio diverges as $g\to\infty$ in this unconstrained task model.
\end{proposition}
\begin{proof}
Substitute the two cost levels into the preceding price index. If $\sigma<1$, $g^{\sigma-1}\to0$ and the remaining positive non-AI weight yields the finite limit. The Cobb--Douglas product gives $P=g^{-a}$. For $\sigma>1$, the bracket tends to infinity with a positive exponent. These are statements about a fixed technology and resource experiment; they do not assert unbounded physical gains when real compute capacity is imposed.
\end{proof}

For example, $a=1/2$, $\sigma=1/2$, $g=4$ and $F/R=1/10$ give $P=9/16$ and $G=8/5$, by direct substitution. The corresponding unlimited-task-efficiency ceiling is $18/5$, even before adding finite task capacities.

\section{Explicit capacity-constrained allocation}
Now impose bounds $0<\bar q_j\leq\infty$ on task outputs. Assume $R-F<\sum_jc_j\bar q_j$, interpreting an infinite bound in the usual way, so not every task can be saturated.

\begin{theorem}[One-dimensional determination under capacity bounds]
There is a unique $K>0$ solving
\[
 \sum_j c_j\min\{\bar q_j,K\omega_jc_j^{-\sigma}\}=R-F.\tag{2}
\]
The unique maximizing allocation is
$q_j^*=\min\{\bar q_j,K\omega_jc_j^{-\sigma}\}$. It can be found by monotone scalar search or by sorting the saturation thresholds $\bar q_jc_j^\sigma/\omega_j$. Its output is no larger than the unconstrained output $Q^*=(R-F)/P(c)$.
\end{theorem}
\begin{proof}
The left side of (2) is continuous, starts at zero and increases strictly until all finite tasks saturate; its supremum exceeds $R-F$. Thus it has exactly one relevant solution. Unsaturated coordinates obey the common derivative-per-cost first-order condition. At a saturated coordinate the cap lies at or below its unconstrained amount at this $K$, so its marginal output per resource is at least the common multiplier. These are exactly the upper-bound multiplier inequalities. Concavity makes them sufficient for a global optimum, and strict curvature along any nonproportional positive direction gives uniqueness; proportional distinct allocations cannot both exhaust the same positive budget. Sorting the breakpoints makes the left side piecewise linear in $K$, so the correct interval and then $K$ can also be solved explicitly. Restricting the feasible set cannot increase output.
\end{proof}

\section{Measured TFP and the remaining identification problem}
\begin{proposition}[Output gains need not equal measured TFP gains]
For a fixed accounting convention $A=Y/(K_m^\alpha L_m^{1-\alpha})$, with positive measured inputs and $0<\alpha<1$, two endpoints satisfy
\[
 \log(A_1/A_0)=\log(Y_1/Y_0)
 -\alpha\log(K_{m,1}/K_{m,0})-(1-\alpha)\log(L_{m,1}/L_{m,0}).
\]
In particular, with unchanged measured labor and output ratio $G>1$, the measured TFP change is negative exactly when $K_{m,1}/K_{m,0}>G^{1/\alpha}$.
\end{proposition}
\begin{proof}
Take logarithms of the declared accounting identity at both endpoints and subtract. The strict sign condition follows by rearrangement. The mapping of compute and organizational investment into $K_m$ is part of the accounting convention and is not supplied by task experiments.
\end{proof}

The calculation requires causal task costs, substitution elasticity, deployment costs, available capacity and an accounting mapping. A task experiment alone does not identify these joint primitives or the finite-horizon investment path. The theorem treats optional deployment but fixes entry and the task set. It makes no inference about future unspecified capabilities or autonomous research acceleration.

\section{Reference}
Erik Brynjolfsson, Anton Korinek and Ajay Agrawal (2024), \emph{The Economics of Transformative AI: A Research Agenda}, October 23 draft, Section 3.1, pp.~9--10. \url{https://digitaleconomy.stanford.edu/app/uploads/2024/11/ETAI-White-Paper.pdf}. The source supplies the growth-and-constraints agenda; the resource-constrained CES benchmark is derived here.

\end{document}
